\SourceLeaf{352}
\[
\left.\begin{aligned}
\gamma&=\frac{1}{\surd{\bigl(1-(\frac{aA}{aa+G})^2-(\frac{bB}{bb+G})^2\bigr)}}\\[9pt]
\gamma'&=\frac{1}{\surd{\bigl((\frac{aA}{aa-G'})^2+(\frac{bB}{bb-G'})^2-1\bigr)}}\\[9pt]
\gamma''&=\frac{1}{\surd{\bigl((\frac{aA}{aa-G''})^2+(\frac{bB}{bb-G''})^2-1\bigr)}}
\end{aligned}\right\}\quad[15]
\]
Ex his aequationibus rite cum 5, 8, 11 combinatis etiam sequitur:
\[
\left.\begin{aligned}
\gamma&=\surd{\frac{G}{(\frac{AG}{aa+G})^2+(\frac{BG}{bb+G})^2+CC}}\\[9pt]
\gamma'&=\surd{\frac{G'}{(\frac{AG'}{aa-G'})^2+(\frac{BG'}{bb-G'})^2+CC}}\\[9pt]
\gamma''&=\surd{\frac{G''}{(\frac{AG''}{aa-G''})^2+(\frac{BG''}{bb-G''})^2+CC}}
\end{aligned}\right\}\quad[16]
\]
Hae posteriores expressiones ostendunt, nullam quantitatum $G,G',G''$ negativam esse posse, siquidem $\gamma,\gamma',\gamma''$ debent esse reales.

In casu itaque eo, ubi non est $C=0$, necessario $G$ aequalis statui debet radici positivae aequationis $B$, patetque adeo, $-G'$ aequalem esse debere
alteri radici negativae, atque $-G''$ aequalem alteri\HistoricalNote{26}; utram vero radicem pro
$-G'$, utram pro $-G''$ adoptemus, prorsus arbitrarium erit.

Quoties $C=0$, punctumque attractum intra curvam situm, duas radices
negativas aequationis 13 necessario pro $-G'$ et $-G''$ adoptare et proin $G=0$
statuere oportet. Quoniam vero in hoc casu formula prima in 16 fit indeterminata, formulam primam in 15 eius loco retinebimus, quae suppeditat
\[
\begin{aligned}
&\gamma=\frac{1}{\surd{\bigl(1-\frac{AA}{aa}-\frac{BB}{bb}\bigr)}}
\end{aligned}
\]
Quoties autem pro $C=0$ punctum attractum extra ellipsin iacet. aequa\SourceLeaf{353}%
tionis 13 radix positiva statuenda est $=G$, atque vel negativa $=-G'$, et
$G''=0$, vel radix negativa $=-G''$, et $G'=0$; coëfficientem $\gamma''$ vel $\gamma'$ vero
inveniemus per formulam
\[
\begin{aligned}
&\frac{1}{\surd{\bigl(\frac{AA}{aa}+\frac{BB}{bb}-1\bigr)}}
\end{aligned}
\]
Ceterum in casu iam excluso, ubi punctum attractum in ipsa circumferentia ellipsis situm supponeretur, coëfficientes $\gamma$ et $\gamma'$, vel $\gamma$ et $\gamma''$ evaderent infiniti,
quod indicat, transformationem nostram ad hunc casum omnino non esse applicabilem.

\AttractionArticle{8}
Quamquam formulae 15, 16 ad determinationem coëfficientium $\gamma,\gamma',\gamma''$
sufficere possent, tamen etiam elegantiores assignare licet. Ad hunc finem multiplicabimus aequationem [5] per $aabb-GG$, unde prodit, levi reductione facta,
\[
\begin{aligned}
&\frac{aaAA(bb+G)}{aa+G}-AAG+\frac{bbBB(aa+G)}{bb+G}-BBG+\frac{aabbCC}{G}-CCG\\[6pt]
&\qquad{}=aabb-GG
\end{aligned}
\]
Sed e natura aequationis cubicae fit
\[
\begin{aligned}
&\text{summa radicum}\quad G-G'-G''=AA+BB+CC-aa-bb\\[8pt]
&\text{productum radicum}\quad GG'G''=aabbCC
\end{aligned}
\]
Hinc aequatio praecedens transit in sequentem:
\[
\begin{aligned}
&\frac{aaAA(bb+G)}{aa+G}+\frac{bbBB(aa+G)}{bb+G}+G'G''-G(G-G'-G''+aa+bb)\\[6pt]
&\qquad{}=aabb-GG
\end{aligned}
\]
quam etiam sic exhibere licet
\[
\begin{aligned}
&\frac{aaAA(bb+G)}{aa+G}+\frac{bbBB(aa+G)}{bb+G}-(aa+G)(bb+G)\\[6pt]
&\qquad{}+(G+G')(G+G'')=0
\end{aligned}
\]
Hinc valor coëfficientis $\gamma$ e formula prima in [15] transmutatur in sequentem:
\[
\begin{aligned}
&\gamma=\surd{\frac{(aa+G)(bb+G)}{(G+G')(G+G'')}}\quad.\,.\,.\,.\,.\,.\,.\,.\quad[17]
\end{aligned}
\]
Per analysin prorsus similem invenitur
\[
\begin{aligned}
&\gamma'=\surd{\frac{(aa-G')(bb-G')}{(G+G')(G''-G')}}\quad.\,.\,.\,.\,.\,.\,.\,.\quad[18]
\end{aligned}
\]
\[
\begin{aligned}
&\gamma''=\surd{\frac{(aa-G'')(bb-G'')}{(G+G'')(G'-G'')}}\quad.\,.\,.\,.\,.\,.\,.\quad[19]
\end{aligned}
\]
\SourceLeaf{354}
\par Postquam coëfficientes $\gamma,\gamma',\gamma''$ inventi sunt, reliqui $\alpha,\beta,\alpha',\beta',\alpha'',\beta''$
inde per formulas 3, 4, 6, 7, 9, 10 derivabuntur.

\AttractionArticle{9}
Signa expressionum radicalium, per quas $\gamma,\gamma',\gamma''$ determinavimus, ad lubitum accipi posse facile perspicitur. Operae autem pretium est, inquirere, quomodo signum quantitatis $\varepsilon$ cum signis istis nexum sit. Ad hunc finem consideremus aequationem tertiam in III art. 3.
\[
\begin{aligned}
&\varepsilon\gamma=\alpha'\beta''-\beta'\alpha''
\end{aligned}
\]
quae per formulas 6, 7, 9, 10 transmutatur in hanc:
\[
\begin{aligned}
&\varepsilon\gamma=\frac{abAB\gamma'\gamma''}{(aa-G')(bb-G'')}-\frac{abAB\gamma'\gamma''}{(aa-G'')(bb-G')}\\[8pt]
&=\frac{ab(aa-bb)AB(G''-G')\gamma'\gamma''}{(aa-G')(aa-G'')(bb-G')(bb-G'')}
\end{aligned}
\]
Sed e consideratione aequationis 13 facile deducimus
\[
\begin{aligned}
&(aa+G)(aa-G')(aa-G'')=aa(aa-bb)AA\\[8pt]
&(bb+G)(bb-G')(bb-G'')=-bb(aa-bb)BB
\end{aligned}
\]
Hinc aequatio praecedens fit
\[
\begin{aligned}
&\varepsilon\gamma=\frac{(aa+G)(bb+G)(G'-G'')\gamma'\gamma''}{ab(aa-bb)AB}
\end{aligned}
\]
quae combinata cum aequatione 17 suppeditat
\[
\begin{aligned}
&\gamma\gamma'\gamma''=\frac{\varepsilon ab(aa-bb)AB}{(G+G')(G+G'')(G'-G'')}
\end{aligned}
\]
Hinc patet, si pro $-G'$ electa sit aequationis cubicae radix negativa absolute maior, simulque coëfficientes $\gamma,\gamma',\gamma''$ omnes positive accepti sint, $\varepsilon$ idem
signum nancisci, quod habet $AB$, idemque evenire, si his quatuor conditionibus, vel omnibus vel duabus ex ipsis, contraria acta sint, oppositum vero, si uni
vel tribus conditionibus adversatus fueris. Ceterum sequentes adhuc relationes
notare convenit, e praecedentibus facile derivandas:
\SourceLeaf{355}
\[
\begin{aligned}
&\alpha\alpha'\alpha''=\frac{\varepsilon aabAAB}{(G+G')(G+G'')(G'-G'')}\\[8pt]
&\beta\beta'\beta''=-\frac{\varepsilon abbABB}{(G+G')(G+G'')(G'-G'')}\\[8pt]
&\alpha\beta=\frac{abAB}{(G+G')(G+G'')}\\[8pt]
&\alpha'\beta'=-\frac{abAB}{(G+G')(G'-G'')}\\[8pt]
&\alpha''\beta''=\frac{abAB}{(G+G'')(G'-G'')}
\end{aligned}
\]

\AttractionArticle{10}

Formulae nostrae quibusdam casibus indeterminatae fieri possunt, quos
seorsim considerare oportet. Ac primo quidem discutiamus casum eum, ubi aequationis
cubicae radices negativae $-G'$, $-G''$ aequales fiunt, unde, per formulas
18, 19, coëfficientes $\gamma'$, $\gamma''$ valores infinitos nancisci videntur, qui autem
revera sunt indeterminati.

Statuendo in formula 14, $g=G$, patet, ut duo valores ipsius $x$, i.\,e. ut
$-G'$ et $-G''$ fiant aequales, necessario esse debere
\[
\begin{aligned}
&AB=0,\quad aa-bb-\frac{aaAA}{aa+G}+\frac{bbBB}{bb+G}=0
\end{aligned}
\]
\par Hinc facile intelligitur, quum $aa-bb$ natura sua sit vel quantitas positiva,
vel $=0$, esse debere
\[
\begin{aligned}
&B=0
\end{aligned}
\]
\[
\begin{aligned}
&aa-bb=\frac{aaAA}{aa+G},\quad\text{sive}\quad aa+G=\frac{aaAA}{aa-bb}
\end{aligned}
\]
Substituendo hos valores in aequatione 14, fit
\[
\begin{aligned}
&G'=G''=bb
\end{aligned}
\]
Substituendo porro valorem $x=-bb$ in aequatione cubica 13, prodit
\[
\begin{aligned}
&(aa-bb)(CC+bb)=bbAA
\end{aligned}
\]
Quoties haec aequatio conditionalis simul cum aequatione $B=0$ locum habet,
casus, quem hic tractamus, adducitur. Et quum fiat
\[
\begin{aligned}
&G=\frac{aaAA}{aa-bb}-aa=\frac{aaCC}{bb}
\end{aligned}
\]
\noindent\SourceLeaf{356}
formula 17 suppeditat
\[
\begin{aligned}
&\gamma=\surd{\frac{aabbAA}{(aa-bb)(aaCC+b^4)}}
=\surd{\frac{aaCC+aabb}{aaCC+b^4}}
\end{aligned}
\]
ac dein formulae 3, 4
\[
\begin{aligned}
&\alpha=\frac{\gamma(aa-bb)}{aA}=\frac{\gamma bbA}{a(CC+bb)}
=\surd{\frac{bb(aa-bb)}{aaCC+b^4}}
\\[6pt]
&\qquad{}=\surd{\frac{b^4AA}{(CC+bb)(aaCC+b^4)}}
\end{aligned}
\]
\[
\begin{aligned}
&\beta=0
\end{aligned}
\]
\par Valores coëfficientium $\gamma'$, $\gamma''$ per formulas 18, 19 in hoc casu indeterminati
manent, atque sic etiam valores coëfficientium reliquorum $\alpha'$, $\beta'$, $\alpha''$, $\beta''$. Nihilominus
per unum horum coëfficientium omnes quinque reliqui exprimi possunt,
e.\,g. fit per formulam 6
\[
\begin{aligned}
&\alpha'=\frac{\gamma'aA}{aa-bb}
\end{aligned}
\]
ac dein
\[
\begin{aligned}
&\beta'=\surd{(1-\alpha'\alpha'+\gamma'\gamma')},\\[6pt]
&
\gamma''=\surd{(\gamma\gamma-1-\gamma'\gamma')},\\[6pt]
&
\alpha''=\frac{\gamma''aA}{aa-bb},\\[6pt]
&
\beta''{}=\surd{(1-\alpha''\alpha''+\gamma''\gamma'')}
\end{aligned}
\]
Sed concinnius hoc ita perficitur. Ex
\[
\begin{aligned}
&\gamma\gamma=1+\alpha\alpha,\quad
\alpha\alpha'=\gamma\gamma',\quad
1=\alpha'\alpha'+\beta'\beta'-\gamma'\gamma'
\end{aligned}
\]
sequitur
\[
\begin{aligned}
&\beta'\beta'+\frac{\gamma'\gamma'}{\alpha\alpha}
=1-\alpha'\alpha'+\frac{\gamma\gamma\gamma'\gamma'}{\alpha\alpha}=1
\end{aligned}
\]
\par Quapropter statuere possumus
\[
\begin{aligned}
&\beta'=\cos f,\quad\gamma'=\alpha\sin f,\quad\alpha'=\gamma\sin f
\end{aligned}
\]
Dein vero e formulis
\[
\begin{aligned}
&\varepsilon\alpha''=\beta\gamma'-\gamma\beta',\quad
\varepsilon\beta''=\gamma\alpha'-\alpha\gamma',\quad
\varepsilon\gamma''=\beta\alpha'-\alpha\beta',\quad
\varepsilon\varepsilon=1
\end{aligned}
\]
invenimus
\[
\begin{aligned}
&\alpha''=-\varepsilon\gamma\cos f,\quad
\beta''=\varepsilon\sin f,\quad
\gamma''=-\varepsilon\alpha\cos f
\end{aligned}
\]
Valor anguli $f$ hic arbitrarius est, nec non pro lubitu statui poterit vel $\varepsilon=+1$
vel $\varepsilon=-1$.

\AttractionArticle{11}

Si $G'$, $G''$ sunt inaequales, valores coëfficientium $\gamma$, $\gamma'$, $\gamma''$ per formulas
17, 18, 19 indeterminati esse nequeunt, sed quoties aliqua quantitatum
\noindent\SourceLeaf{357}
$aa-G'$, $bb-G'$, $aa-G''$, $bb-G''$ evanescit, valor coëfficientis $\alpha'$, $\beta'$, $\alpha''$, $\gamma''$
per formulam 6, 7, 9, 10 resp. indeterminatus manere primo aspectu videtur,
quod tamen secus se habere levis attentio docebit.

Supponamus e.\,g., esse $aa-G'=0$, fietque, per aequationem 18, $\gamma'=0$,
nec non per aequationem 7, $\beta'=0$ (siquidem non fuerit simul $aa=bb$)
unde necessario esse debet $\alpha'=\pm1$. Si vero simul $aa=bb$, formula, quae
praecedit sextam in art. 5, suppeditat $\alpha'A+\beta'B=0$, quae aequatio cum
$\alpha'\alpha'+\beta'\beta'=1$ iuncta, producit
\[
\begin{aligned}
&\alpha'=\frac{B}{\surd{(AA+BB)}},\quad
\beta'=\frac{-A}{\surd{(AA+BB)}}
\end{aligned}
\]
\par Hae expressiones manifesto indeterminatae esse nequeunt, nisi simul fuerit
$A=0$, $B=0$; tunc vero ad casum in art. praec. iam consideratum delabemur.

\AttractionArticle{12}

Postquam duodecim quantitates $G$, $G'$, $G''$, $\alpha$, $\alpha'$, $\alpha''$, $\beta$. $\beta'$, $\beta''$, $\gamma$, $\gamma'$, $\gamma''$ complete
determinare docuimus, ad evolutionem differentialis $\mathrm dE$ progredimur.
Statuamus
\[
\begin{aligned}
&t=\gamma+\gamma'\cos T+\gamma''\sin T\quad
\mathinner{\ldotp\;\ldotp\;\ldotp\;\ldotp\;\ldotp\;\ldotp\;\ldotp\;\ldotp}\quad[20]
\end{aligned}
\]
ita ut fiat
\[
\begin{aligned}
&t\cos E=\alpha+\alpha'\cos T+\alpha''\sin T\quad
\mathinner{\ldotp\;\ldotp\;\ldotp\;\ldotp\;\ldotp\;\ldotp\;\ldotp}\quad[21]
\end{aligned}
\]
\[
\begin{aligned}
&t\sin E=\beta+\beta'\cos T+\beta''\sin T\quad
\mathinner{\ldotp\;\ldotp\;\ldotp\;\ldotp\;\ldotp\;\ldotp\;\ldotp}\quad[22]
\end{aligned}
\]
Hinc deducimus
\[
\begin{aligned}
&t\,\mathrm dE=\cos E\,\mathrm d.\,t\sin E-\sin E\,\mathrm d.\,t\cos E\\[8pt]
&=\cos E(\beta''\cos T-\beta'\sin T)\mathrm dT
-\sin E(\alpha''\cos T-\alpha'\sin T)\mathrm dT
\end{aligned}
\]
adeoque
\[
\begin{aligned}
&tt\,\mathrm dE=(\alpha\beta''-\alpha''\beta)\cos T\,\mathrm dT
+(\alpha'\beta-\beta'\alpha)\sin T\,\mathrm dT
\\[6pt]
&\qquad{}+(\alpha'\beta''-\beta'\alpha'')\mathrm dT\\[8pt]
&=\varepsilon\gamma'\cos T\,\mathrm dT
+\varepsilon\gamma''\sin T\,\mathrm dT
+\varepsilon\gamma\,\mathrm dT=\varepsilon t\,\mathrm dT
\end{aligned}
\]
sive
\[
\begin{aligned}
&t\,\mathrm dE{}=\varepsilon\,\mathrm dT\quad
\mathinner{\ldotp\;\ldotp\;\ldotp\;\ldotp\;\ldotp\;\ldotp\;\ldotp\;\ldotp\;\ldotp\;\ldotp\;\ldotp\;\ldotp\;\ldotp\;\ldotp\;\ldotp\;\ldotp\;\ldotp\;\ldotp\;\ldotp\;\ldotp\;\ldotp}\quad[23]
\end{aligned}
\]
\par Observare convenit, quantitatem $t$ natura sua semper positivam esse, si
coëfficiens $\gamma$ sit positivus, vel semper negativam, si $\gamma$ sit negativus. Quum enim
sit $(\gamma'\cos T+\gamma''\sin T)^2+(\gamma''\cos T-\gamma'\sin T)^2=\gamma'\gamma'+\gamma''\gamma''=\gamma\gamma-1$, erit semper
\noindent\SourceLeaf{358}
$\gamma'\cos T+\gamma''\sin T$, sine respectu signi, minor quam $\gamma$. Hinc concludimus, quoties
$\varepsilon\gamma$ sit quantitas positiva, variables $E$ et $T$ semper simul crescere; quoties
autem $\varepsilon\gamma$ sit quantitas negativa, necessario alteram variabilem semper decrescere,
dum altera augeatur.

\AttractionArticle{13}

Nexus inter variables $E$ et $T$ adhuc melius illustratur per ratiocinia sequentia.
Statuendo $\surd{(\gamma\gamma-1)}=\delta$, ita ut fiat $\delta\delta=\alpha\alpha+\beta\beta=\gamma'\gamma'+\gamma''\gamma''$, ex
aequationibus 20, 21, 22 deducimus
\[
\begin{aligned}
&t(\delta+\alpha\cos E+\beta\sin E)\\[8pt]
&\quad=\gamma\delta+\alpha\alpha+\beta\beta
+(\gamma'\delta+\alpha\alpha'+\beta\beta')\cos T
+(\gamma''\delta+\alpha\alpha''+\beta\beta'')\sin T\\[8pt]
&\quad=(\gamma+\delta)(\delta+\gamma'\cos T+\gamma''\sin T)
\end{aligned}
\]
\par Perinde ex aequationibus 21, 22 sequitur
\[
\begin{aligned}
&t(\alpha\sin E-\beta\cos E)=\varepsilon(\gamma'\sin T-\gamma''\cos T)
\end{aligned}
\]
Hae aequationes, statuendo
\[
\begin{aligned}
&\frac{\alpha}{\delta}=\cos L,\quad
\frac{\beta}{\delta}=\sin L,\quad
\frac{\gamma'}{\delta}=\cos M,\quad
\frac{\gamma''}{\delta}=\sin M
\end{aligned}
\]
nanciscuntur formam sequentem:
\[
\begin{aligned}
&t(1+\cos(E-L))=(\gamma+\delta)(1+\cos(T-M))
\end{aligned}
\]
\[
\begin{aligned}
&t\sin(E-L)=\varepsilon\sin(T-M)
\end{aligned}
\]
unde fit per divisionem, propter $(\gamma+\delta)(\gamma-\delta)=1$,
\[
\begin{aligned}
&\mathrm{tang}\,\frac12(E-L)=\varepsilon(\gamma-\delta)\mathrm{tang}\,\frac12(T-M)
\end{aligned}
\]
\[
\begin{aligned}
&\mathrm{tang}\,\frac12(T-M)=\varepsilon(\gamma+\delta)\mathrm{tang}\,\frac12(E-L)
\end{aligned}
\]
\par Hinc non solum eadem conclusio derivatur, ad quam in fine art. praec. deducti
sumus, sed insuper etiam patet, si valor ipsius $E$ crescat 360 gradibus, valorem
ipsius $T$ tantundem vel crescere vel diminui, prout $\varepsilon\gamma$ sit vel quantitas
positiva vel negativa. Ceterum statuendo $\delta=\mathrm{tang}\,N$, $\gamma=\sec N$, manifesto erit
\[
\begin{aligned}
&\gamma-\delta=\mathrm{tang}(45^0-\frac12N),\quad
\gamma+\delta=\mathrm{tang}(45^0+\frac12N)
\end{aligned}
\]
\SourceLeaf{359}
\AttractionArticle{14}

E combinatione aequationum 20, 21, 22 cum aequationibus art. 5 obtinemus:
\[
\begin{aligned}
&at(A-a\cos E)=\alpha G-\alpha'G'\cos T-\alpha''G''\sin T\\[8pt]
&bt(B-b\sin E)=\beta G-\beta'G'\cos T-\beta''G''\sin T
\end{aligned}
\]
\par Statuendo itaque brevitatis gratia
\[
\begin{aligned}
&(\alpha G-\alpha'G'\cos T-\alpha''G''\sin T)\\[8pt]
&\quad(\gamma-e\alpha+(\gamma'-e\alpha')\cos T+(\gamma''-e\alpha'')\sin T)=aX\\[8pt]
&(\beta G-\beta'G'\cos T-\beta''G''\sin T)\\[8pt]
&\quad(\gamma-e\alpha+(\gamma'-e\alpha')\cos T+(\gamma''-e\alpha'')\sin T)=bY\\[8pt]
&C(\gamma+\gamma'\cos T+\gamma''\sin T)\\[8pt]
&\quad(\gamma-e\alpha+(\gamma'-e\alpha')\cos T+(\gamma''-e\alpha'')\sin T)=Z
\end{aligned}
\]
fit
\[
\begin{aligned}
&\mathrm d\xi=\frac{\varepsilon X\,\mathrm dT}{2\pi t^3\rho^3},\qquad
\mathrm d\eta=\frac{\varepsilon Y\,\mathrm dT}{2\pi t^3\rho^3},\qquad
\mathrm d\zeta=\frac{\varepsilon Z\,\mathrm dT}{2\pi t^3\rho^3}
\end{aligned}
\]
\par Sed habetur
\[
\begin{aligned}
&t\rho=\pm\surd{(G+G'\cos T^2+G''\sin T^2)}
\end{aligned}
\]
signo superiore vel inferiore valente, prout $t$ est quantitas positiva vel negativa
($\rho$ enim natura sua semper positive accipitur), i. e. prout coëfficiens $\gamma$ est positivus vel negativus. Hinc
\[
\begin{aligned}
&\frac{\varepsilon\,\mathrm dT}{2\pi t^3\rho^3}
=\pm\frac{\mathrm dT}{2\pi(G+G'\cos T^2+G''\sin T^2)^{\frac32}}
\end{aligned}
\]
ubi signum ambiguum a signo quantitatis $\gamma\varepsilon$ pendet.

Ut iam valores ipsarum $\xi,\eta,\zeta$ obtineamus, integrationes differentialium exsequi oportet,
a valore ipsius $T$, cui respondet $E=0$, usque ad valorem, cui respondet $E=360^0$,
sive etiam (quod manifesto eodem redit) a valore ipsius $T$, cui respondet valor arbitrarius ipsius $E$,
usque ad valorem, cui respondet valor ipsius $E$ auctus $360^0$;
licebit itaque integrare a $T=0$ usque ad $T=360^0$, quoties $\varepsilon\gamma$ est quantitas positiva,
vel a $T=360^0$ usque ad $T=0$, quoties $\varepsilon\gamma$ est negativa.
Manifesto itaque, independenter a signo ipsius $\varepsilon\gamma$, erit:
\[
\begin{aligned}
&\xi=\int\frac{X\,\mathrm dT}{2\pi(G+G'\cos T^2+G''\sin T^2)^{\frac32}}\\[8pt]
&\eta=\int\frac{Y\,\mathrm dT}{2\pi(G+G'\cos T^2+G''\sin T^2)^{\frac32}}\\[8pt]
&\zeta=\int\frac{Z\,\mathrm dT}{2\pi(G+G'\cos T^2+G''\sin T^2)^{\frac32}}
\end{aligned}
\]
integrationibus a $T=0$ usque ad $T=360^0$ extensis.
\SourceLeaf{360}
\AttractionArticle{15}

Nullo negotio perspicitur, integralia
\[
\begin{aligned}
&\int\frac{\cos T\,\mathrm dT}{(G+G'\cos T^2+G''\sin T^2)^{\frac32}}\\[8pt]
&\int\frac{\sin T\,\mathrm dT}{(G+G'\cos T^2+G''\sin T^2)^{\frac32}}\\[8pt]
&\int\frac{\cos T\sin T\,\mathrm dT}{(G+G'\cos T^2+G''\sin T^2)^{\frac32}}
\end{aligned}
\]
a $T=180^0$ usque ad $T=360^0$ extensa obtinere valores aequales iis, quos nanciscuntur,
si a $T=0$ usque ad $T=180^0$ extendantur, sed signis oppositis affectos;
quapropter ista integralia a $T=0$ usque ad $T=360^0$ extensa manifesto fiunt $=0$.
Hinc colligimus. esse
\[
\begin{aligned}
&\xi{}=\int\frac{\begin{gathered}((\gamma-e\alpha)\alpha G-(\gamma'-e\alpha')\alpha'G'\cos T^2\\-(\gamma''-e\alpha'')\alpha''G''\sin T^2)\,\mathrm dT\end{gathered}}{2\pi a(G+G'\cos T^2+G''\sin T^2)^{\frac32}}\\[8pt]
&\eta{}=\int\frac{\begin{gathered}((\gamma-e\alpha)\beta G-(\gamma'-e\alpha')\beta'G'\cos T^2\\-(\gamma''-e\alpha'')\beta''G''\sin T^2)\,\mathrm dT\end{gathered}}{2\pi b(G+G'\cos T^2+G''\sin T^2)^{\frac32}}\\[8pt]
&\zeta{}=\int\frac{\begin{gathered}((\gamma-e\alpha)\gamma+(\gamma'-e\alpha')\gamma'\cos T^2\\+(\gamma''-e\alpha'')\gamma''\sin T^2)C\,\mathrm dT\end{gathered}}{2\pi(G+G'\cos T^2+G''\sin T^2)^{\frac32}}
\end{aligned}
\]
integralibus a $T=0$ usque ad $T=360^0$ extensis.
Quodsi itaque valores integralium, eadem extensione acceptorum,
\[
\begin{aligned}
&\int\frac{\cos T^2\,\mathrm dT}{2\pi((G+G')\cos T^2+(G+G'')\sin T^2)^{\frac32}}\\[8pt]
&\int\frac{\sin T^2\,\mathrm dT}{2\pi((G+G')\cos T^2+(G+G'')\sin T^2)^{\frac32}}
\end{aligned}
\]
per $P,Q$ denotamus, erit
\[
\begin{aligned}
&a\xi=((\gamma-e\alpha)\alpha G-(\gamma'-e\alpha')\alpha'G')P\\[6pt]
&\qquad{}+((\gamma-e\alpha)\alpha G-(\gamma''-e\alpha'')\alpha''G'')Q\\[8pt]
&b\eta=((\gamma-e\alpha)\beta G-(\gamma'-e\alpha')\beta'G')P\\[6pt]
&\qquad{}+((\gamma-e\alpha)\beta G-(\gamma''-e\alpha'')\beta''G'')Q\\[8pt]
&\zeta=((\gamma-e\alpha)\gamma+(\gamma'-e\alpha')\gamma')CP\\[6pt]
&\qquad{}+((\gamma-e\alpha)\gamma+(\gamma''-e\alpha'')\gamma'')CQ
\end{aligned}
\]
quo pacto problema nostrum complete solutum est.
\SourceLeaf{361}
\AttractionArticle{16}

Quod attinet ad quantitates $P,Q$, manifesto quidem utraque fit
\[
\begin{aligned}
&=\frac{1}{2(G+G')^{\frac32}}
\end{aligned}
\]
quoties $G'=G''$, in omnibus vero reliquis casibus ad transscendentes sunt referendae.
Quas quomodo per series exprimere liceat, abunde constat.
Lectoribus autem gratum fore speramus, si hacce occasione determinationem harum aliarumque transscendentium
per algorithmum peculiarem expeditissimum explicemus, quo per multos iam abhinc annos frequenter usi sumus,
et de quo alio loco copiosius agere propositum est.

Sint $m,n$ duae quantitates positivae, statuamusque
\[
\begin{aligned}
&m'=\frac12(m+n),\qquad n'=\surd{mn}
\end{aligned}
\]
ita ut $m',n'$ resp. sit medium arithmeticum et geometricum inter $m$ et $n$.
Medium geometricum semper positive accipi supponemus. Perinde fiat
\[
\begin{aligned}
&m''=\frac12(m'+n'),\qquad n''=\surd{m'n'}\\[8pt]
&m'''=\frac12(m''+n''),\qquad n'''=\surd{m''n''}
\end{aligned}
\]
et sic porro, quo pacto series $m,m',m'',m'''$ etc., atque $n,n',n'',n'''$ etc.
versus \emph{limitem communem} rapidissime convergent, quem per $\mu$ designabimus,
atque simpliciter \emph{medium arithmetico-geometricum} inter $m$ et $n$ vocabimus.
Iam demonstrabimus, $\frac1\mu$ esse valorem integralis
\[
\begin{aligned}
&\int\frac{\mathrm dT}{2\pi\surd{(mm\cos T^2+nn\sin T^2)}}
\end{aligned}
\]
a $T=0$ usque ad $T=360^0$ extensi.

\emph{Demonstr.} Supponamus, variabilem $T$ ita per aliam $T'$ exprimi, ut fiat
\[
\begin{aligned}
&\sin T=\frac{2m\sin T'}{(m+n)\cos T^{\prime 2}+2m\sin T^{\prime 2}}
\end{aligned}
\]
perspicieturque facile, dum $T'$ a valore $0$ usque ad $90^0,180^0,270^0,360^0$ augeatur,
etiam $T$ (etsi inaequalibus intervallis) a $0$ usque ad $90^0,180^0,270^0,360^0$ crescere.
Evolutione autem rite facta, invenitur esse
\[
\begin{aligned}
&\frac{\mathrm dT}{\surd{(mm\cos T^2+nn\sin T^2)}}
=\frac{\mathrm dT'}{\surd{(m'm'\cos T^{\prime 2}+n'n'\sin T^{\prime 2})}}
\end{aligned}
\]
